%%============================================================
%% JSAI 2026 Poster
%% geDIG: Gauge what Knowledge Graph needs.
%%
%% Build: lualatex poster.tex (run twice for cross-references)
%%
%% Layout: NeurIPS-inspired 2-column hybrid, A0 portrait
%% - Hero title + QR
%% - When 問題 (60%) + Hero 数値 (40%)
%% - Figure 1 hybrid (full width)
%% - Formulation (50%) + entropy_sign (50%)
%% - Spatial KG (50%) + Semantic KG w/ bar chart (50%)
%% - Discussion SP→β₁ (50%) + OCR (50%)
%% - Limitations + Links (full width)
%%============================================================

\documentclass[a0paper, portrait, 25pt, margin=1.2cm]{tikzposter}

%% ---- 日本語 (lualatex + luatexja) ----
\usepackage{luatexja}
\usepackage{graphicx}
\usepackage{hyperref}
\usepackage{xcolor}

%% ---- 数式・表 ----
\usepackage{amsmath, amssymb}
\usepackage{booktabs}
\usepackage{array}

%% ---- TikZ ----
\usepackage{tikz}
\usetikzlibrary{arrows.meta, positioning, shapes.geometric, fit, calc, backgrounds}

%% ---- コマンド定義 ----
\newcommand{\F}{\mathcal{F}}
\newcommand{\gednorm}{\Delta \mathrm{EPC}}
\newcommand{\ignorm}{\Delta H}
\newcommand{\spnorm}{\Delta \mathrm{SP}}
\newcommand{\btnorm}{\Delta \beta_1}

%% ---- カスタムカラー ----
\definecolor{heroblue}{HTML}{2563EB}
\definecolor{herored}{HTML}{DC2626}
\definecolor{herogreen}{HTML}{059669}
\definecolor{softblue}{HTML}{DBEAFE}
\definecolor{softred}{HTML}{FEE2E2}
\definecolor{softgreen}{HTML}{D1FAE5}

%% ---- ポスターテーマ ----
\tikzposterlatexaffectionproofoff
\usetheme{Simple}
\usecolorstyle{Britain}
\useblockstyle{Default}
\usetitlestyle{Filled}

%% ---- タイトル設定 ----
\title{\parbox{0.92\linewidth}{\centering
  動的知識グラフの探索・統合を制御する統一ゲージの提案\\[4mm]
  \normalsize \textit{Gauge what Knowledge Graph needs.}
}}
\author{宮内 和義 \quad (Kazuyoshi Miyauchi)}
\institute{Independent Researcher \quad
  \texttt{miyauchikazuyoshi@gmail.com} \quad
  \href{https://github.com/miyauchikazuyoshi/InsightSpike-AI}
       {\texttt{github.com/miyauchikazuyoshi/InsightSpike-AI}}}
\titlegraphic{}

%%============================================================
\begin{document}
%%============================================================

\maketitle

%%============================================================
%% Row 1: When 問題 (全幅)
%%============================================================
\block{1. 本研究の問い: When 問題}{
\begin{minipage}[t]{0.62\linewidth}
\vspace{0mm}
\textbf{動的知識グラフにおける根本問題}:
エージェントが新しい情報に出会ったとき,
「\textbf{いつ受け入れ, いつ探索を続けるべきか?}」

\medskip
\textbf{従来}: KL / ELBO / FEP / MDL 等の\textbf{確率的評価}.
構造項と確率項を分離せず, 位相的変化を暗黙的にしか扱えない.

\medskip
\textbf{本研究}: 構造コストと情報利得を\textbf{独立項としてスカラー化}し,
単一ゲージ \(\F\) で定量化.
\end{minipage}%
\hfill
\begin{minipage}[t]{0.35\linewidth}
\vspace{0mm}
\begin{center}
  \Large \(\F = \gednorm - \lambda \bigl( s \cdot \ignorm + \gamma \cdot \spnorm \bigr)\)
\end{center}
\medskip
\textbf{2 段ゲート AG/DG} により \(\F\) の符号で
\textbf{\textcolor{heroblue}{受け入れ}} / 保留 /
\textbf{\textcolor{herored}{拒否}} の 3 判定.

\medskip
\textbf{貢献 3 点}: (1) SP 定義改善 (CLS 集中度), (2) entropy\_sign \(s\) 導入,
(3) 介入実験による因果検証.
\end{minipage}
}

%%============================================================
%% Row 1.5: Hero 数値 (全幅、大カード 3 枚)
%%============================================================
\block[bodyinnersep=8mm, titleoffsety=0cm]{}{
\vspace{-5mm}
\begin{center}
\begin{tikzpicture}
  \node[fill=softblue, draw=heroblue, line width=1.5pt,
        rounded corners=10pt, text width=10.5cm, align=center,
        inner sep=8pt] (cardA) at (0, 0) {
    {\color{heroblue}\veryHuge\textbf{98\%}}\\[2mm]
    {\Large \textbf{Maze} ゴール到達率}\\[1mm]
    {\normalsize $15 \times 15$ grid, 250 steps}
  };
  \node[fill=softred, draw=herored, line width=1.5pt,
        rounded corners=10pt, text width=10.5cm, align=center,
        inner sep=8pt] (cardB) at (13, 0) {
    {\color{herored}\veryHuge\textbf{4.6\%}}\\[2mm]
    {\Large \textbf{Transformer} 介入差}\\[1mm]
    {\normalsize Positive $-$ Negative (因果的実証)}
  };
  \node[fill=softgreen, draw=herogreen, line width=1.5pt,
        rounded corners=10pt, text width=10.5cm, align=center,
        inner sep=8pt] (cardC) at (26, 0) {
    {\color{herogreen}\veryHuge\textbf{18K}}\\[2mm]
    {\Large \textbf{OCR} params $\to$ 73.53\%}\\[1mm]
    {\normalsize 外部プロジェクトでの実証}
  };
\end{tikzpicture}
\end{center}
}

%%============================================================
%% Row 2: Figure 1 hybrid (full width) + caption
%%============================================================
\block{2. Why three terms? \textemdash{} AG/DG 二段ゲートと 3 判定}{
\begin{minipage}[t]{0.38\linewidth}
\vspace{0mm}
\begin{center}
  \resizebox{0.96\linewidth}{!}{%%
%% Figure 1 Hybrid: AG/DG 二段ゲートと When 問題の 3 判定
%% orbit (空間フレーム) + matchstick (3 ケース) の統合版
%%
%% 外側円 = Multi-hop Reach (DG 領域、ΔSP で評価)
%% 内側円 = 0-hop Orbit (AG 領域、ΔEPC, ΔH で評価)
%% 3 ケース = Case A (Insight) / B (力仕事) / C (崩壊)
%%

\begin{tikzpicture}[scale=0.85, every node/.style={transform shape}]

  %% ============ 背景の 2 つの円 ============

  % Multi-hop Reach (DG 領域): 外側円、薄い赤背景
  \fill[red!5] (0,0) circle (6.5cm);
  \draw[red!60, thick] (0,0) circle (6.5cm);

  % 0-hop Orbit (AG 領域): 内側円、薄い青背景
  \fill[blue!8] (0,0) circle (3.0cm);
  \draw[blue!60, thick, dashed] (0,0) circle (3.0cm);

  %% ============ 領域ラベル ============

  % Multi-hop Reach: 外側円の上部、円の外
  \node[red!70, font=\small\bfseries] at (0, 7.0) {Multi-hop Reach (DG)};
  \node[red!70, font=\footnotesize] at (0, 6.55) {評価: $\Delta \mathrm{SP}$};

  % 0-hop Orbit: 内側円の左外側（Case A との重なりを避ける）
  \node[blue!70, font=\small\bfseries, anchor=east] at (-3.2, 1.6)
    {0-hop Orbit (AG)};
  \node[blue!70, font=\footnotesize, anchor=east] at (-3.2, 1.0)
    {評価: $\Delta \mathrm{EPC}, \Delta H$};

  %% ============ Query (中心) ============

  \node[star, star points=5, fill=red, draw=red!80, inner sep=4pt,
        label={[font=\small\bfseries, red!80]below:Query $Q$}] (Q) at (0,0) {};

  %% ============ Case A: Insight (Multi-hop 圏で赤太線) ============

  % ノード
  \node[circle, draw=black, fill=white, inner sep=2.5pt] (A1) at (-4.8, 3.8) {};
  \node[circle, draw=black!60, fill=white, inner sep=1.5pt] (A2) at (-5.5, 4.2) {};
  \node[circle, draw=black!60, fill=white, inner sep=1.5pt] (A3) at (-4.5, 4.8) {};
  \draw[black!60] (A1) -- (A2);
  \draw[black!60] (A1) -- (A3);

  % 赤太線で Multi-hop 短絡
  \draw[red, line width=2.5pt, -{Latex[length=4mm]}] (A1) -- (Q);

  % ラベル (ノードの上方、0-hop ラベルと重ならない位置)
  \node[align=center, font=\small, fill=white, draw=red!40, rounded corners,
        inner sep=3pt] at (-5.2, 5.5)
       {\textbf{Case A: Insight}\\
        $\Delta \mathrm{SP}{=}+1,\ F{<}0$\\
        \textcolor{red!80}{\textbf{DG 発火 → 受け入れ}}};

  %% ============ Case B: 力仕事 (0-hop 圏内で編集) ============

  % 0-hop 圏内のノード
  \node[circle, draw=black, fill=white, inner sep=2.5pt] (B1) at (1.4, 1.3) {};
  \node[circle, draw=black, fill=white, inner sep=2.5pt] (B2) at (1.4, -0.3) {};

  % 普通の黒エッジ (Q に届かない編集)
  \draw[black, line width=1.2pt] (B1) -- (B2);

  % ラベル (0-hop 円の右外側、Multi-hop 圏内の下部)
  \node[align=center, font=\small, fill=white, draw=black!30, rounded corners,
        inner sep=3pt] at (4.8, -1.7)
       {\textbf{Case B: 力仕事}\\
        $\Delta \mathrm{SP}{=}0,\ F{\approx}0$\\
        \textcolor{black!60}{\textbf{発火なし → 保留}}};

  %% ============ Case C: 崩壊 (Multi-hop 既存短絡の消失) ============

  % 外側のノード (元々 Q と短絡していた)
  \node[circle, draw=black, fill=white, inner sep=2.5pt] (C1) at (4.5, -4.5) {};
  \node[circle, draw=black!60, fill=white, inner sep=1.5pt] at (5.3, -4.0) {};
  \node[circle, draw=black!60, fill=white, inner sep=1.5pt] at (5.0, -5.3) {};

  % 薄い点線で元の短絡を示す (削除されたことを表現)
  \draw[red!30, line width=2.5pt, dashed] (C1) -- (Q);

  % 削除マーク (×)
  \node[red, font=\LARGE\bfseries] at (2.3, -2.3) {\texttimes};

  % ラベル (外側円の内側、下部)
  \node[align=center, font=\small, fill=white, draw=red!40, rounded corners,
        inner sep=3pt] at (5.0, -5.3)
       {\textbf{Case C: 崩壊}\\
        $\Delta \mathrm{SP}{=}-1,\ F{>}0$\\
        \textcolor{red!80}{\textbf{逆発火 → 拒否}}};

  %% ============ 凡例 (共通) ============

  \node[align=left, font=\scriptsize, fill=white, draw=black!20,
        rounded corners, inner sep=4pt] at (-5.5, -5.5)
       {\textbf{凡例}\\
        \textcolor{red}{\rule[0.5mm]{5mm}{1.2mm}} 有効短絡 $(\Delta \mathrm{SP}{\gg}0)$\\
        \textcolor{red!40}{\rule[0.5mm]{5mm}{0.6mm}}$\,\cdots$ 消失した短絡\\
        \rule[0.5mm]{5mm}{0.4mm} 通常の編集};

\end{tikzpicture}
}
\end{center}
\small\centering
\textbf{Figure 1}: AG/DG 二段ゲートと 3 判定の空間的構造
\end{minipage}%
\hfill
\begin{minipage}[t]{0.60\linewidth}
\vspace{0mm}
\textbf{2 段ゲートの必然性}\\
Query $Q$ (Transformer の CLS, 迷路のゴール) を中心に,
geDIG は 2 段階で新情報を評価する.

\smallskip
\begin{itemize}\setlength\itemsep{1mm}
  \item \textcolor{heroblue}{\textbf{AG (内側円, 0-hop)}}: 近傍の曖昧性を検知.
    $g_0 = \gednorm - \lambda \cdot \ignorm$ の大小で発火.
  \item \textcolor{herored}{\textbf{DG (外側円, multi-hop)}}: 遠方への短絡が有効か確認.
    \(\spnorm \gg 0\) なら発火.
\end{itemize}
\small AG だけでは全探索になりコスト爆発, DG だけでは新発見不能.
2 段組み合わせで\textbf{曖昧性検知 $\to$ 有効性検証}の流れを作る.

\medskip
\textbf{3 ケースの判定 (Figure 1 参照)}

\smallskip
\begin{itemize}\setlength\itemsep{1mm}
  \item \textbf{Case A: Insight} \; $\spnorm{=}+1$, $F{<}0$\\
    Multi-hop 圏に新短絡が張られる. CLS への新経路発見.\\
    $\to$ DG 発火 $\to$ \textbf{\textcolor{heroblue}{受け入れ}}
    (\S 5B の Positive 介入と対応)
  \item \textbf{Case B: 力仕事} \; $\spnorm{=}0$, $F{\approx}0$\\
    0-hop 圏内の編集. CLS 集中度に寄与しない.\\
    $\to$ 発火なし $\to$ \textbf{保留}
  \item \textbf{Case C: 崩壊} \; $\spnorm{=}-1$, $F{>}0$\\
    既存の Multi-hop 短絡が消失. CLS への有効経路断絶.\\
    $\to$ 逆発火 $\to$ \textbf{\textcolor{herored}{拒否}}
    (\S 5B の Negative 介入と対応)
\end{itemize}

\medskip
\textbf{なぜ位相項 $\spnorm$ が必要か}\\
Case A/B/C は同じ編集コスト (\(\gednorm{=}1\))
かつ類似の情報量変化 (\(\ignorm\)) でも, \textbf{CLS への経路集中度
\(\spnorm\)} が $+1/0/-1$ に分岐する. 確率的評価 (KL 等) は
$\spnorm$ の情報を見逃す.

\medskip
$\to$ 位相項を独立に F に組み込むことで,
\textbf{質的に異なる 3 判定}が初めて工学的に実現される.
\textbf{これが geDIG の When 問題への解法.}
\end{minipage}
}

%%============================================================
%% Row 3: Formulation (50%) + entropy_sign (50%)
%%============================================================
\begin{columns}

\column{0.5}
\block{3. Formulation: \(\F\) の定義}{
\begin{center}
\(\F = \gednorm - \lambda \bigl( s \cdot \ignorm + \gamma \cdot \spnorm \bigr)\)
\end{center}

\begin{center}
\begin{tabular}{@{}ll@{}}
\toprule
項 & 意味 \\
\midrule
\(\gednorm\) & 構造変更コスト (編集距離) \\
\(\ignorm\) & 情報量変化 (Shannon) \\
\(\spnorm\) & \textbf{Shortcut Purity 変化}\\
             & (CLS への経路集中度) \\
\(\lambda, \gamma\) & 重み係数 (分位点校正) \\
\(s\) & entropy\_sign (\(\pm 1\)) \\
\bottomrule
\end{tabular}
\end{center}

\medskip
\(\F < 0\): \textbf{\textcolor{heroblue}{受け入れ}} (コストに見合う更新)\\
\(\F \approx 0\): 保留 \\
\(\F > 0\): \textbf{\textcolor{herored}{拒否}}
}

\column{0.5}
\block{4. entropy\_sign \(s\): 学習フェーズ適応}{
学習フェーズによって \(\ignorm\) の望ましい方向が異なる:

\medskip
\begin{center}
\begin{tabular}{@{}lcl@{}}
\toprule
フェーズ & \(s\) & 解釈 \\
\midrule
\textbf{構築} (迷路, 事前学習) & \textbf{\color{heroblue}+1} & 延伸 = 利得 \\
\textbf{特化} (Fine-tuning) & \textbf{\color{herored}-1} & 集中 = 利得 \\
\bottomrule
\end{tabular}
\end{center}

\medskip
これにより,
\begin{itemize}\setlength\itemsep{1mm}
  \item \textcolor{heroblue}{\textbf{空間 KG}} (迷路, 構築)
  \item \textcolor{herored}{\textbf{意味 KG}} (Transformer, 特化)
\end{itemize}
を同一ゲージで評価できる.

\medskip
\textbf{単一数式 \(\F\) と \(s\) の組合せで 2 ドメイン統一}.
}

\end{columns}

%%============================================================
%% Row 4: Experiments — Spatial KG (50%) + Semantic KG (50%)
%%============================================================
\begin{columns}

\column{0.5}
\block{5A. Spatial KG: 迷路ナビゲーション}{
\textbf{設定}: $15 \times 15$ grid, 部分観測, 250 steps,
\(s=+1\) (構築フェーズ).

\medskip
\textbf{結果}:
\begin{center}
\begin{tabular}{@{}lc@{}}
\toprule
指標 & 値 \\
\midrule
ゴール到達率 & \textbf{\color{heroblue}98\%} \\
候補エッジ拒否率 & \textbf{\color{heroblue}98\%} \\
AG/DG 発火 & デッドエンド/ループ閉鎖と一致 \\
\bottomrule
\end{tabular}
\end{center}

\medskip
\textbf{観察}: AG はデッドエンドで即座に発火, DG はループ閉鎖時に圧縮を実行.
\(\F\) の定性的な動作が理論予測と\textbf{一致}.

\medskip
\textit{\small 注: 単一設定 ($15 \times 15$) の結果. 3-seed 以上の統計的検証は Future Work.}
}

\column{0.5}
\block{5B. Semantic KG: Transformer 介入実験}{
\textbf{設定}: DistilBERT / SST-2, \(s=-1\) (特化フェーズ),
Positive / Negative 介入で \(\F\) の方向を操作.

\medskip
\textbf{介入実験結果} (accuracy):
\begin{center}
\begin{tikzpicture}[scale=0.9]
  \def\barh{0.55}

  % Positive (upper bar, blue, F down = accept)
  \fill[heroblue!60, draw=heroblue] (0, 2*\barh) rectangle (8.56, 3*\barh);
  \node[anchor=east, heroblue!80!black] at (-0.15, 2.5*\barh)
    {\small Positive (\(\F{\downarrow}\))};
  \node[anchor=west] at (8.6, 2.5*\barh) {\textbf{85.6\%}};

  % Baseline (middle, gray)
  \fill[gray!40, draw=gray!70] (0, \barh) rectangle (8.68, 2*\barh);
  \node[anchor=east] at (-0.15, 1.5*\barh) {\small Baseline};
  \node[anchor=west] at (8.6, 1.5*\barh) {\textbf{86.8\%}};

  % Negative (lower bar, red, F up = reject)
  \fill[herored!50, draw=herored] (0, 0) rectangle (8.10, \barh);
  \node[anchor=east, herored!80!black] at (-0.15, 0.5*\barh)
    {\small Negative (\(\F{\uparrow}\))};
  \node[anchor=west] at (8.6, 0.5*\barh) {\textbf{81.0\%}};

  % X軸
  \draw[->, gray!70] (0, -0.15) -- (10, -0.15);
  \node[gray!80] at (5, -0.55) {\scriptsize Accuracy (\%)};

  % 差分ブラケット
  \draw[thick, herored!80!black, decorate,
        decoration={brace, amplitude=4pt, mirror}]
    (8.3, 3*\barh+0.05) -- (8.3, 0-0.05)
    node[midway, right, xshift=5mm, align=left]
    {\scriptsize \textbf{4.6\%}\\\scriptsize 差};
\end{tikzpicture}
\end{center}

\medskip
\textbf{Positive − Negative = 4.6\% の精度差}\\
$\to$ \(\F\) は\textbf{因果的}な学習シグナル (観測でなく介入で実証).

\medskip
\textit{\small 注: 単一 seed の結果. 3-seed 再現 + 他モデル (BERT-base) 拡張を計画.}
}

\end{columns}

%%============================================================
%% Row 5: Discussion — SP→β₁ (50%) + OCR (50%)
%%============================================================
\begin{columns}

\column{0.5}
\block{6. Discussion: SP \(\to \beta_1\) への一般化}{
SP はクエリ $Q$ (CLS) への経路集中度 = \textbf{クエリ依存の位相項}.

\medskip
\textit{採択論文後の自然な拡張}:
$Q$ 依存性を外した一般化として\textbf{第一 Betti 数} $\beta_1$
(独立サイクル数) が考えられる.

\medskip
$\beta_1$ の性質:
\begin{itemize}\setlength\itemsep{1mm}
  \item \textbf{位相不変量} (重み非依存)
  \item \textbf{離散量} (\(\btnorm \in \mathbb{Z}\))
  \item \textbf{次元フリー} (計算量 \(O(V+E)\))
\end{itemize}

\medskip
並行記録実験で Spearman \(r > 0.7\) (SP, $\beta_1$ 間)
$\to$ SP は $\beta_1$ の特殊ケースと解釈可能.

\medskip
Figure 1 の外側円 (Multi-hop Reach) は $Q$ 依存のクエリ圏. $\beta_1$
ではこのクエリ依存性を外して閉路の数そのものを位相量とする.
}

\column{0.5}
\block{7. 外部実証: 視覚処理 (OCR)}{
外部プロジェクト
\href{https://github.com/miyauchikazuyoshi/vector-based-cnn-ocr}
     {\texttt{vector-based-cnn-ocr}}
で geDIG 原理を検証 (Font $\to$ 手書き cross-domain).

\medskip
\begin{center}
\begin{tabular}{@{}lc@{}}
\toprule
指標 & 値 \\
\midrule
パラメータ数 & \textbf{\color{herogreen}18K} (標準 CNN の $1/100+$) \\
精度 & \textbf{\color{herogreen}73.53\%} \\
DG 情報効率 / AG & \textbf{\color{herogreen}5倍} \\
\bottomrule
\end{tabular}
\end{center}

\medskip
\textbf{示唆}:
\begin{itemize}\setlength\itemsep{1mm}
  \item 構造を先に与えれば params は桁違いに減少
  \item \textbf{DG (確定情報) > AG (探索情報)} の効率性
  \item 空間/意味 KG 以外でも geDIG が機能
\end{itemize}

\medskip
\textit{\small 詳細: \texttt{docs/research/gedig\_core\_theory\_unified.md \S 8.2}}
}

\end{columns}

%%============================================================
%% Row 6: Limitations + Conclusion + Links (full width)
%%============================================================
\begin{columns}

\column{0.60}
\block{8. Limitations (誠実な範囲明示)}{
\small
\begin{itemize}\setlength\itemsep{1mm}
  \item \textbf{統計的検定}: 迷路/Transformer ともに単一 seed.
    3-seed 以上の再現と paired comparison を計画.
  \item \textbf{モデル規模}: DistilBERT のみ. BERT-base / GPT-2 / Pythia 拡張は Future Work.
  \item \textbf{学習初期化}: 事前学習済みモデルでの介入のみ.
    ランダム初期化は Pythia で計画中.
  \item \textbf{\(\beta_1\) への完全移行}: SP-\(\beta_1\) 相関 $r > 0.7$ は確認済だが,
    Transformer での \(\beta_1\) 直接検証は継続中.
  \item \textbf{\(\lambda, \gamma\) 感度分析}: 分位点校正による経験的選択のみ.
\end{itemize}
}

\column{0.40}
\block{9. Future Work \& Links}{
\small
\textbf{次の一手}:
\begin{itemize}\setlength\itemsep{1mm}
  \item Pythia 学習過程全体の観察
  \item SP $\to \beta_1$ 完全移行 (Transformer)
  \item Grokking 相転移での \(\F\) + curl 観測
\end{itemize}

\medskip
\textbf{Resources}:\\
\href{https://github.com/miyauchikazuyoshi/InsightSpike-AI}
     {\textcolor{heroblue}{\small \texttt{github.com/miyauchikazuyoshi/InsightSpike-AI}}}\\
\href{https://doi.org/10.5281/zenodo.19454110}
     {\textcolor{heroblue}{\small DOI: \texttt{10.5281/zenodo.19454110}}}\\
\texttt{miyauchikazuyoshi@gmail.com}
}

\end{columns}

\end{document}
